\documentclass[
  final,
  screen,
  editorial=false,
  informations=false,
  membership-reminder=false,
  detailedtoc=2,
  number=42,
  date=2026-08
]{letgut}

\addbibresource{references.bib}
\title{La lettre des outils ouverts}
\author{Association GUTenberg}

\begin{document}
\maketitle
% letgut's end-of-document bookmark hook assumes article's author command is
% still defined after its custom title page has run.
\makeatletter
\ProvideDocumentCommand{\author}{m}{\gdef\@author{#1}}
\makeatother

\inputarticle*{editorial}

\section{Une chaîne éditoriale reproductible}

Un bon bulletin technique doit être lisible à l'écran comme sur papier. La
classe \texttt{letgut} rassemble une mise en page de bulletin, des marges
éditoriales et une bibliographie locale afin que chaque numéro puisse être
reconstruit sans accès réseau. Cette démonstration cite une source de référence
sur la composition numérique~\cite{knuth1984}.

\begin{announcement}[color=orange!65!black,colback=yellow!10]{À retenir}
  Gardez les données éditoriales, les images et le fichier bibliographique dans
  le même dossier que le numéro distribué.
\end{announcement}

\subsection{Un exemple de code}

\begin{lstlisting}[language={[LaTeX]TeX},caption={Une section minimale}]
\section{Titre}
Un paragraphe de contenu, suivi de \cite{lamport1994}.
\printbibliography
\end{lstlisting}

\printbibliography[heading=subtitle,title={Références}]
\end{document}
